\begin{afterword}
\summaryhead

The post extends ``Entangled Truths, Contagious Lies,'' from two days earlier. A lie is out of place in a lawful, connected world, so protecting it requires further lies: first about connected facts, then about the general laws that link them, then about the rules of reasoning. Creationists, the author says, lie about evolutionary theory and then about what ``theory'' means.

The same happens in self-deception. In an imagined exchange, a believer in a dragon, asked for evidence, denies that beliefs require evidence, then denies the reasons for that rule, then retreats to a ``separate magisterium.'' A single ``Lie That Must Be Protected'' can block progress in rationality. So wherever lies are protected, bad general epistemology, ``the Dark Side,'' comes into being. ``Everyone has a right to their own opinion'' is offered as a possible example. Most who repeat such ideas are duped rather than deceivers, and the idea is not to be used to dismiss arguments without refuting them. The post ends by asking readers to list more examples.

\responsehead

The post's model is coherent, and parts of it are old and sound. That a belief can be kept against any evidence by revising other beliefs, even the rules of logic, is Quine's point of 1951, though the post does not mention Quine. The creationist arguments it describes are recognizable. The post also limits itself in ways worth noting: it corrects its own pebble example in a footnote, says most people who repeat bad epistemology are ``more duped than duplicitous,'' and forbids using its idea in place of refutation.

The empirical claims have little support. That bad rules of thinking arise in this way, and that a single protected belief ``can block someone's progress into advanced rationality,'' are supported by the model, by the phrase ``But how else?'', and by a made-up exchange with a believer in a dragon. No person whose reasoning was blocked is named. The one case, creationism, is described in general terms: no creationist is quoted, and no rule of reasoning is traced, with evidence, to a lie it was made to protect.

The one example of a Dark Side proverb shows the problem. The post asks where ``Everyone has a right to their own opinion'' came from, and does not look. In its political sense it is a right to hold opinions without interference, a claim about coercion, not about evidence. The question also judges the proverb by the purpose of whoever first said it, and three paragraphs later the post states the rule that an idea must be refuted ``for itself,'' not by its inventor's intentions. The closing request, to list more memes ``spawned by the Dark Side,'' asks readers for more classifications of the same kind.

Two smaller points. The post says ``most creationists'' lie about evolution, and later that most repeaters of bad epistemology are duped; if ``lie'' has its ordinary sense, the two claims are hard to reconcile. And its list of what a biologist would see includes proteins held by ``weak van der Waals forces,'' with no reason given why weak bonds count against design.

In short: a coherent model of how protecting one belief corrupts others, supported by a made-up exchange rather than cases, and illustrated by a proverb that the post judges by its supposed origin, contrary to its own rule that ideas be refuted for themselves.
\end{afterword}
